\documentclass[11pt,a4paper]{article}

\usepackage[utf8]{inputenc}
\usepackage[T1]{fontenc}
\usepackage[spanish]{babel}
\usepackage[a4paper,margin=2.3cm]{geometry}
\usepackage{amsmath,amssymb,amsfonts,amsthm}
\usepackage{mathtools}
\usepackage{bm}
\usepackage{booktabs}
\usepackage{enumitem}
\usepackage{xcolor}
\usepackage{hyperref}
\usepackage{tcolorbox}

\hypersetup{
  colorlinks=true,
  linkcolor=blue!50!black,
  urlcolor=blue!60!black,
  citecolor=blue!40!black
}

\newcommand{\BHF}{B_{\mathrm{hf}}}
\newcommand{\dEQ}{\Delta E_{Q}}
\newcommand{\code}[1]{\texttt{\small #1}}
\newcommand{\file}[1]{\textcolor{black!70}{\texttt{\small #1}}}
\DeclareMathOperator{\diag}{diag}
\DeclareMathOperator{\tr}{tr}
\DeclareMathOperator{\erfc}{erfc}
\DeclareMathOperator*{\argmin}{arg\,min}
\DeclareMathOperator*{\argmax}{arg\,max}

\newtcolorbox{nota}[1][Nota]{
  colback=blue!4, colframe=blue!45!black, fonttitle=\bfseries,
  title={#1}, left=4pt, right=4pt, top=3pt, bottom=3pt
}

\title{\textbf{Manual matemático del motor de ajuste \emph{Mossbauer}}\\[3pt]
\large Modelo directo, ajuste discreto y ajuste de distribución}
\author{Departamento de Química Física · UAM}
\date{\today}

\begin{document}
\maketitle

\begin{abstract}\noindent
Este manual documenta, de forma autocontenida, la matemática completa del motor de
ajuste de espectros Mössbauer de $^{57}$Fe implementado en la aplicación. Cubre
(i)~el \emph{modelo directo} —perfiles de línea, posiciones e intensidades de las
seis líneas del sextete con tratamiento de primer orden y de Kündig del
acoplamiento magnético–cuadrupolar, dobletes y singletes—; (ii)~el \emph{ajuste
discreto} —estadística de los datos en modos Gauss y Poisson, propagación de la
incertidumbre de calibración, funcional de coste con pérdidas robustas,
optimización por región de confianza con arranque múltiple y opción global,
covarianza, diagnósticos y bootstrap—; y (iii)~el \emph{ajuste de distribución}
$P(\BHF)$ o $P(\dEQ)$ —discretización, núcleo, regularización de Tikhonov y de
variación total, condicionamiento, grados de libertad efectivos, selección
automática del parámetro de regularización, barras de error y distribuciones
paramétricas. Se incluyen referencias a fichero y función del código.
\end{abstract}

\tableofcontents
\bigskip

%======================================================================
\part{Modelo directo}
%======================================================================

\section{Espectro de transmisión}\label{sec:fwd}

Tras plegar los datos brutos en torno al punto de doblado, el espectro queda
definido sobre $N$ canales con velocidades $\{v_i\}_{i=1}^{N}$ equiespaciadas en
$[-v_{\max}, v_{\max}]$ y transmisión normalizada $\{y_i\}$, con $y_i\approx 1$
fuera de los picos de absorción. El modelo directo es una superposición lineal de
absorciones sobre una línea base:
\begin{equation}
  \boxed{\;
  T(v;\bm\theta) \;=\; b + s\,v \;-\; \sum_{c=1}^{N_c} A_c(v;\bm\theta_c)
  \;}
  \label{eq:Tmodel}
\end{equation}
donde $b$ es la línea base (multiplicativa, los datos están normalizados por el
factor \code{norm\_factor}), $s$ la pendiente, y cada $A_c$ es la absorción de un
componente (sextete, doblete o singlete). Es un modelo de \emph{absorbente
delgado}: lineal en la profundidad de cada componente. Implementado en
\file{core/physics.py}, función \code{total\_model}.

\section{Perfiles de línea}\label{sec:lineshape}

Cada línea elemental usa un perfil normalizado a altura unidad en su centro. La
aplicación ofrece dos perfiles, seleccionables globalmente.

\subsection{Lorentziana}
\begin{equation}
  L(v;v_0,\gamma) \;=\; \frac{\gamma^{2}}{(v-v_0)^{2}+\gamma^{2}},
  \label{eq:lorentz}
\end{equation}
con $\gamma$ la semianchura a media altura (HWHM). Cumple $L(v_0)=1$.

\subsection{Pseudo-Voigt (Faddeeva)}
La convolución de una Lorentziana (HWHM $\gamma$) con una Gaussiana (anchura
$\sigma$) es el perfil de Voigt, evaluado mediante la función de Faddeeva
$w(z)=e^{-z^2}\erfc(-iz)$:
\begin{equation}
  \widetilde V(v;v_0,\gamma,\sigma) \;=\;
   \frac{1}{\sigma\sqrt{2\pi}}\,\Re\!\left[w\!\left(\frac{(v-v_0)+i\gamma}{\sigma\sqrt2}\right)\right].
\end{equation}
Para que el perfil tenga altura exactamente $1$ en $v_0$ \emph{independientemente
del muestreo}, se normaliza por su valor analítico en el pico:
\begin{equation}
  V(v;v_0,\gamma,\sigma) \;=\; \frac{\widetilde V(v;v_0,\gamma,\sigma)}{\widetilde V(v_0;v_0,\gamma,\sigma)},
  \qquad
  \widetilde V(v_0) = \frac{1}{\sigma\sqrt{2\pi}}\,\Re\!\left[w\!\left(\tfrac{i\gamma}{\sigma\sqrt2}\right)\right].
  \label{eq:voigt}
\end{equation}
Como $w(iy)=e^{y^2}\erfc(y)\in\mathbb R^{+}$ para $y\in\mathbb R$, el denominador
es real y positivo. En el límite $\sigma\to0$, $V\to L$.

\begin{nota}[Normalización al pico]
Normalizar por el máximo \emph{discreto} de la malla (como hacían versiones
previas) subestima el pico real cuando ningún $v_i$ cae en $v_0$, inflando la
amplitud y haciendo que la profundidad dependa del muestreo. La forma
\eqref{eq:voigt} elimina ese sesgo.
\end{nota}

\section{Sextete magnético}\label{sec:sextet}

Un sextete surge del desdoblamiento Zeeman nuclear combinado, en su caso, con la
interacción cuadrupolar. Las seis líneas corresponden a las transiciones
permitidas ($\Delta m=0,\pm1$) entre el estado fundamental $I_g=\tfrac12$ y el
excitado $I_e=\tfrac32$. La absorción es
\begin{equation}
  A_{\text{sext}}(v) \;=\; d\,\sum_{j=1}^{6} W_j\, P\!\left(v;v_{0,j},\gamma_j\right),
  \label{eq:Asext}
\end{equation}
con $d$ la profundidad, $P$ el perfil (\S\ref{sec:lineshape}), $W_j$ las
intensidades (\S\ref{sec:intensities}) y $\gamma_j$ las anchuras
\begin{equation}
  \bm\gamma = \gamma_1\,[\,1,\;\gamma_2,\;\gamma_3,\;\gamma_3,\;\gamma_2,\;1\,],
\end{equation}
donde $\gamma_2,\gamma_3$ son anchuras \emph{relativas}. Las posiciones
$v_{0,j}$ se obtienen por uno de dos tratamientos.

\subsection{Tratamiento de primer orden}
El cuadrupolo se trata como una perturbación rígida sobre el patrón Zeeman:
\begin{equation}
  v_{0,j} \;=\; \frac{\BHF}{B_0}\,r_j^{(0)} \;+\; \delta \;+\; s_j\,\dEQ,
  \qquad j=1,\dots,6,
  \label{eq:sext_pos1}
\end{equation}
con $B_0=33\,$T el campo de referencia, $\delta$ el desplazamiento isomérico,
$r^{(0)}$ las posiciones calibradas a $B_0$ (Apéndice~\ref{ap:calib}) y el patrón
cuadrupolar
\begin{equation}
  \bm s = \left[\,+\tfrac12,\,-\tfrac12,\,-\tfrac12,\,-\tfrac12,\,-\tfrac12,\,+\tfrac12\,\right].
\end{equation}
Válido cuando $\dEQ \ll \mu_N g_e \BHF$ y el eje del gradiente de campo eléctrico
(EFG) es paralelo a $\vec B$.

\subsection{Tratamiento de Kündig (Hamiltoniano completo axial)}\label{sec:kundig}
Para EFG axial (parámetro de asimetría $\eta=0$) y ángulo $\beta$ arbitrario entre
$\vec B$ y el eje principal $V_{zz}$, las posiciones se obtienen diagonalizando el
Hamiltoniano del estado excitado. En la base $\{\lvert\tfrac32\rangle,
\lvert\tfrac12\rangle, \lvert{-}\tfrac12\rangle, \lvert{-}\tfrac32\rangle\}$:
\begin{equation}
  \mathcal H_e \;=\; \omega_e\, I_z \;+\; \frac{\dEQ}{6}\bigl(3 I_{z'}^2 - I(I+1)\bigr),
  \qquad I_{z'} = I_z\cos\beta + I_x\sin\beta,
  \label{eq:Hkundig}
\end{equation}
con $I(I+1)=\tfrac{15}{4}$ y $\omega_e=(\BHF/B_0)\,\omega_e^{(0)}$ el
desdoblamiento Zeeman del excitado. La diagonalización numérica de la matriz
$4\times4$ da los autovalores $E^{e}_{m}$ y autovectores. El fundamental
($I_g=\tfrac12$) es Zeeman puro:
\begin{equation}
  E^{g}_{\pm 1/2} = \pm\tfrac12\,\omega_g, \qquad \omega_g=(\BHF/B_0)\,\omega_g^{(0)}.
\end{equation}
Las seis posiciones son las diferencias de energía de las transiciones permitidas,
asignando a cada autovalor del excitado su $m_e$ dominante (por solapamiento del
autovector):
\begin{equation}
\begin{aligned}
  v_{0,1}&=E^{e}_{-3/2}-E^{g}_{-1/2}, & v_{0,2}&=E^{e}_{-1/2}-E^{g}_{-1/2}, & v_{0,3}&=E^{e}_{+1/2}-E^{g}_{-1/2},\\
  v_{0,4}&=E^{e}_{-1/2}-E^{g}_{+1/2}, & v_{0,5}&=E^{e}_{+1/2}-E^{g}_{+1/2}, & v_{0,6}&=E^{e}_{+3/2}-E^{g}_{+1/2},
\end{aligned}
\end{equation}
más $\delta$. En el límite $\beta=0$ reproduce \eqref{eq:sext_pos1} a precisión de
máquina. Implementado en \file{core/hamiltonian.py}.

\subsection{Promedio policristalino}\label{sec:powder}
En una muestra de polvo todas las orientaciones $\beta$ son equiprobables. La
absorción se promedia sobre la esfera:
\begin{equation}
  A_{\text{sext}}^{\text{pol}}(v) \;=\; \frac12\int_0^{\pi}\!\sin\beta\;A_{\text{sext}}(v;\beta)\,d\beta
  \;\approx\; \sum_{k=1}^{N_q} w_k\, A_{\text{sext}}(v;\beta_k),
  \label{eq:powder}
\end{equation}
con nodos y pesos de cuadratura de Gauss–Legendre sobre $x=\cos\beta\in[-1,1]$,
$\beta_k=\arccos x_k$, $w_k=\tfrac12 w_k^{\mathrm{GL}}$ (de modo que
$\sum_k w_k=1$). El cálculo está vectorizado: para $N_q$ orientaciones se
construye y diagonaliza un tensor $(N_q,4,4)$ de una vez. Valor por defecto
$N_q=20$ (error $<0{,}01\%$).

\section{Intensidades y textura}\label{sec:intensities}

Las intensidades de las seis líneas respetan la simetría $1\!\leftrightarrow\!6$,
$2\!\leftrightarrow\!5$, $3\!\leftrightarrow\!4$:
\begin{equation}
  \bm W = [\,I_1,\,I_2,\,I_3,\,I_3,\,I_2,\,I_1\,],\qquad
  I_1=i_3 i_1,\;\; I_2=i_3 i_2,\;\; I_3=i_3.
\end{equation}
La aplicación ofrece dos parametrizaciones por sextete.

\subsection{Intensidades libres}
$i_1,i_2,i_3$ son parámetros independientes del ajuste. No impone la relación
$3{:}2{:}1$; útil para flexibilidad numérica e interpretación empírica.

\subsection{Parámetro de textura}\label{sec:texture}
Para una muestra con ángulo $\theta$ entre el momento magnético y la dirección del
fotón $\gamma$, las amplitudes de dipolo magnético son
\begin{equation}
  W_{1,6}\propto 3(1+\cos^2\theta),\quad W_{2,5}\propto 4\sin^2\theta,\quad W_{3,4}\propto (1+\cos^2\theta).
\end{equation}
Definiendo el parámetro de textura $t=\sin^2\theta\in[0,1]$ y normalizando por
$(1+\cos^2\theta)=2-t$:
\begin{equation}
  \boxed{\;i_1 = 3,\qquad i_2 = \frac{4t}{2-t},\qquad i_3 = 1.\;}
  \label{eq:texture}
\end{equation}
Casos: $t=\tfrac23\Rightarrow 3{:}2{:}1$ (polvo aleatorio); $t=1\Rightarrow
3{:}4{:}1$ (textura plana, lámina $\perp$ haz); $t=0\Rightarrow 3{:}0{:}1$ (eje
fácil $\parallel$ haz). En este modo $t$ es el único parámetro libre de
intensidad; $i_1,i_2,i_3$ se derivan de \eqref{eq:texture}.

\section{Doblete y singlete}

\begin{align}
  A_{\text{dob}}(v) &= d\bigl[\,i_1\,P(v;\delta-\tfrac{\dEQ}{2},\gamma_1) + i_1 i_2\,P(v;\delta+\tfrac{\dEQ}{2},\gamma_1\gamma_2)\,\bigr],\\
  A_{\text{sing}}(v) &= d\,i_1\,P(v;\delta,\gamma_1).
\end{align}

%======================================================================
\part{Ajuste discreto}
%======================================================================

\section{Parametrización y restricciones}\label{sec:disc-param}

El vector de parámetros libres $\bm x\in\mathbb R^{p}$ se compone de los globales
$\{b,s\}$ (más opcionalmente $v_{\max}$ y el centro de doblado $c_{\text{fold}}$) y,
por cada componente activo, un subconjunto de
$\bm\theta_c=(\delta,\dEQ,\BHF,\gamma_1,\gamma_2,\gamma_3,d,i_1,i_2,i_3,t,\beta)$
según su tipo y modo:
\begin{itemize}[leftmargin=*]
  \item \emph{Tipo}: singlete usa $\{\delta,\gamma_1,d,i_1\}$; doblete añade
        $\{\dEQ,\gamma_2,i_2\}$; sextete usa todos.
  \item \emph{Intensidades}: en modo textura, $\{i_1,i_2,i_3\}$ se sustituyen por
        el único parámetro $t$.
  \item \emph{Cuadrupolo}: en modo Kündig con $\beta$ fijo, $\beta$ entra como
        parámetro libre; en modo policristal o primer orden, no.
  \item \emph{Fijos}: parámetros marcados se retiran del vector activo.
  \item \emph{Restringidos lineales} $k_t=\alpha\,k_s+\beta_0$: se aplican en cada
        evaluación (hasta seis iteraciones para cadenas de dependencias).
\end{itemize}
Cada parámetro tiene cotas duras $[\ell_k,u_k]$ (Apéndice~\ref{ap:bounds}).

\section{Estadística de los datos}\label{sec:weighting}

\subsection{Modo Gauss}
El doblado promedia dos canales Poisson independientes; para conteos $c_i$ por
canal plegado, $\mathrm{Var}\approx c_i/2$, de donde la incertidumbre $1\sigma$
del dato normalizado:
\begin{equation}
  \sigma_i \;=\; \frac{\max\!\bigl(\sqrt{c_i/2},\,1\bigr)}{f_{\text{norm}}},
  \qquad \sigma_i\leftarrow\max(\sigma_i,10^{-9}).
  \label{eq:sigma-gauss}
\end{equation}

\subsection{Modo Poisson (IRLS)}
La verosimilitud Poisson se aproxima por mínimos cuadrados reponderados: la
incertidumbre se evalúa sobre el \emph{modelo} (cuentas predichas
$\hat c_i = T_i f_{\text{norm}}$) y se recalcula en cada evaluación del
optimizador, lo que equivale a un paso de Newton sobre $-2\log\mathcal L$ Poisson:
\begin{equation}
  \sigma_i^{\text{Poiss}} \;=\; \frac{\sqrt{\max(\hat c_i/2,\,1)}}{f_{\text{norm}}}
  \;=\; \frac{\sqrt{\max(T_i f_{\text{norm}}/2,\,1)}}{f_{\text{norm}}}.
  \label{eq:sigma-poisson}
\end{equation}

\subsection{Propagación de la incertidumbre de calibración}\label{sec:calib}
Si la calibración de velocidad aporta una incertidumbre $\sigma_{v_{\max}}$, su
efecto sobre el modelo se suma en cuadratura. Como $v_i=v_{\max}(2i/(N-1)-1)$, se
tiene $\partial v_i/\partial v_{\max}=v_i/v_{\max}$, y por la regla de la cadena
\begin{equation}
  \sigma_i^{\text{eff}\,2} \;=\; \sigma_i^2 \;+\;
  \left(\frac{\partial T}{\partial v}\bigg|_i \frac{v_i}{v_{\max}}\,\sigma_{v_{\max}}\right)^{\!2}.
  \label{eq:calib}
\end{equation}
El término pesa más en los flancos de los picos (donde $\partial T/\partial v$ es
grande), que es donde un error de escala sesga $\delta$ y $\BHF$.

\section{Funcional de coste y pérdidas robustas}\label{sec:cost}

El residuo ponderado es
\begin{equation}
  r_i(\bm x) \;=\; \frac{T(v_i;\bm x)-y_i}{\sigma_i^{\text{eff}}},
  \label{eq:residuo}
\end{equation}
y el funcional minimizado
\begin{equation}
  \mathcal C(\bm x) \;=\; \tfrac12 \sum_{i=1}^{N}\rho\!\left(r_i(\bm x)^2\right),
  \label{eq:cost}
\end{equation}
con $\rho$ la función de pérdida: \emph{lineal} $\rho(z)=z$ (mínimos cuadrados),
\emph{soft-}$\ell_1$ $\rho(z)=2(\sqrt{1+z}-1)$, o \emph{Huber}. Las robustas usan
escala $f_{\text{scale}}=3$ (umbral $\sim3\sigma$), reduciendo el peso de canales
con picos espurios. Para $\rho$ lineal, $\mathcal C=\tfrac12\chi^2$.

\section{Optimización}\label{sec:opt}

\subsection{Región de confianza (TRF)}
Se minimiza con \emph{Trust Region Reflective} sobre las cotas:
\begin{equation}
  \hat{\bm x} \;=\; \argmin_{\bm\ell\le\bm x\le\bm u}\; \mathcal C(\bm x),
\end{equation}
resolviendo en cada iteración el subproblema cuadrático con Jacobiano por
diferencias finitas:
\begin{equation}
  \min_{\bm d}\;\tfrac12\|J_k\bm d + \bm r_k\|^{2}
  \;\;\text{s.a.}\;\; \|\bm D_k\bm d\|\le\Delta_k,\;\; \bm\ell-\bm x_k\le\bm d\le\bm u-\bm x_k.
\end{equation}

\subsection{Arranque múltiple}
Se prueban el punto del usuario más 8 reinicios gausianos:
\begin{equation}
  \bm x^{(s)}=\bm x^{(0)}+\bm\xi^{(s)},\quad \xi^{(s)}_k\sim\mathcal N(0,\rho_k^2(u_k-\ell_k)^2),
\end{equation}
con $\rho_k=0{,}12$ para $\{\delta,\dEQ,\BHF,\gamma_1,d,v_{\max}\}$ y $0{,}08$ en el
resto; semilla fija. Se retiene el de menor coste.

\subsection{Optimización global (opcional)}
Como pre-pasada opcional se ejecuta \emph{differential evolution} sobre el coste
escalar $\tfrac12\|\bm r\|^2$ (inicialización Sobol, sin pulido), y su mejor punto
se añade como semilla para el pulido TRF. Cubre mínimos locales por permutación de
etiquetas e intercambios $\BHF_1\!\leftrightarrow\!\BHF_2$ en espectros
multi-sextete que el arranque gausiano no alcanza.

\section{Incertidumbres y diagnósticos}\label{sec:disc-stats}

\subsection{Covarianza}
A partir del Jacobiano final $J=U\Sigma V^\top$ (SVD, valores singulares truncados):
\begin{equation}
  \boxed{\;\widehat{\mathrm{Cov}}(\hat{\bm x}) = V\,\Sigma^{-2}V^\top\cdot\frac{2\,\mathcal C(\hat{\bm x})}{N-p},
  \qquad \hat\sigma_{x_k}=\sqrt{\widehat{\mathrm{Cov}}_{kk}}.\;}
  \label{eq:cov}
\end{equation}
Se reportan correlaciones $|\rho_{kl}|\ge0{,}95$ como aviso de degeneración.

\subsection{Estadísticos de bondad de ajuste}
\begin{align}
  \chi^2 &= \sum_i r_i^2, & \nu&=N-p, & \chi^2_\nu&=\chi^2/\nu,\\
  \mathrm{AIC}&=N\ln(\mathrm{RSS}/N)+2p, & \mathrm{BIC}&=N\ln(\mathrm{RSS}/N)+p\ln N.
\end{align}

\subsection{Diagnósticos del residuo}
Para detectar estructura no aleatoria (con $\tilde r=r-\bar r$):
\begin{align}
  \text{autocorrelación lag-1:}\quad &\rho_1=\frac{\sum_i\tilde r_i\tilde r_{i+1}}{\sum_i\tilde r_i^2},\\
  \text{rachas (Wald–Wolfowitz):}\quad &Z=\frac{R-\mu_R}{\sigma_R},\\
  \text{correlación antisimétrica:}\quad &c_{\text{as}}=\frac{\tilde{\bm r}\cdot(-\tilde{\bm r}^{\text{rev}})}{\tilde{\bm r}\cdot\tilde{\bm r}},
\end{align}
con avisos si $|\rho_1|>0{,}35$, $|Z|>2$ o $c_{\text{as}}>0{,}45$ (esta última
detecta errores de centrado/calibración).

\subsection{Bootstrap}
Estimación Monte Carlo de errores: $M$ réplicas de datos sintéticos
$\tilde y^{(m)}=T(\cdot;\hat{\bm x})+\eta^{(m)}$, con $\eta^{(m)}$ gausiano de
desviación $\sigma_i$ (o remuestreo Poisson real
$\tilde c\sim\mathrm{Poisson}(\hat c)$ en modo Poisson); se reajusta cada réplica y
$\hat\sigma_{x_k}^{\text{boot}}=\mathrm{std}_m(\hat x_k^{(m)})$.

%======================================================================
\part{Ajuste de distribución}
%======================================================================

\section{Discretización y núcleo}\label{sec:dist-kernel}

Se modela el espectro como una densidad continua $P(\BHF)$ (o $P(\dEQ)$)
discretizada en $n$ \emph{bins} de centros equiespaciados
$c_k\in[p_{\min},p_{\max}]$. El \emph{núcleo}
\begin{equation}
  K_{ik} \;=\; A_{\text{sext}}^{(d=1)}\bigl(v_i;\;\BHF=c_k\bigr)
\end{equation}
es la absorción a $v_i$ de un sextete de profundidad unitaria con campo $c_k$
(variable $\BHF$) o con $\dEQ=c_k$ y $\BHF$ fijo (variable $\dEQ$). Las
intensidades del núcleo siguen $3{:}2{:}1$ por construcción.

\section{Sistema lineal aumentado}\label{sec:dist-system}

El forward es
\begin{equation}
  \hat y_i \;=\; b + s\,v_i \;-\; \sum_{k}K_{ik}\,P_k \;-\;\sum_{m}K^{\text{sharp}}_{im}q_m,
\end{equation}
con $P_k\ge0$ los pesos de la distribución y $q_m\ge0$ amplitudes de sextetes
``nítidos'' opcionales. Definiendo la matriz de diseño
\begin{equation}
  X = \bigl[\;\bm 1\;\big|\;\bm v\;\big|\;{-}K\;\big|\;{-}K^{\text{sharp}}\;\bigr],
  \qquad \bm z=(b,s,\bm P^\top,\bm q^\top)^\top,
\end{equation}
y los pesos $W=\diag(1/\sigma_i^2)$, el problema es
\begin{equation}
  \boxed{\;
  \min_{\bm z}\;\;\bigl\|W^{1/2}(X\bm z-\bm y)\bigr\|_2^2 \;+\; \alpha\,\mathcal R(\bm P)
  \quad\text{s.a.}\quad b\ge0,\;\bm P\ge0,\;\bm q\ge0.\;}
  \label{eq:dist-min}
\end{equation}
La pendiente $s$ no tiene cota. Se resuelve con mínimos cuadrados acotados
(\code{lsq\_linear}, reflective Newton con LSMR).

\subsection{Preacondicionamiento por escala de columnas}\label{sec:cond}
Para $\BHF$ pequeño el sextete colapsa y las columnas de $K$ se vuelven casi
colineales con normas muy dispares, lo que mal-condiciona \eqref{eq:dist-min}.
Se reescala cada columna de la distribución $K_k\to K_k/s_k$ con
$s_k=\lVert W^{1/2}K_k\rVert$ y $\tilde P_k=s_k P_k$; el penalizador y los pesos
se ajustan en consecuencia y se deshace la escala al final. Es una
reparametrización \emph{exacta} (solución y grados de libertad invariantes) que
sólo mejora la estabilidad numérica.

\section{Regularización}\label{sec:reg}

\subsection{Tikhonov (suavidad, L2)}
$\mathcal R(\bm P)=\lVert L\,\bm P\rVert_2^2$ con $L$ el operador de \emph{segunda
diferencia} $(n-2)\times n$:
\begin{equation}
  (L\bm P)_k = P_{k-1}-2P_k+P_{k+1},
\end{equation}
que penaliza la curvatura y produce distribuciones suaves. El sistema aumentado
apila $\bigl[\begin{smallmatrix}W^{1/2}X\\ \sqrt\alpha\,L_{\text{ext}}\end{smallmatrix}\bigr]\bm z = \bigl[\begin{smallmatrix}W^{1/2}\bm y\\ \bm 0\end{smallmatrix}\bigr]$.

\subsection{Variación total (bordes, L1)}
$\mathcal R(\bm P)=\lVert D\,\bm P\rVert_1$ con $D$ la \emph{primera diferencia}
$(D\bm P)_k=P_{k+1}-P_k$. Favorece soluciones constantes-a-trozos que preservan
bordes (menos emborronamiento de picos). Se resuelve por mínimos cuadrados
reponderados (IRLS): partiendo de la mayorización
$|t|\le t^2/(2|t^{\text{prev}}|)+\text{cte}$, cada iteración resuelve un Tikhonov
ponderado con
\begin{equation}
  w_k = \frac{1}{\sqrt{(D\bm P^{\text{prev}})_k^2+\varepsilon^2}},
  \qquad \mathcal R\approx \sum_k w_k\,(D\bm P)_k^2,
\end{equation}
reutilizando \code{lsq\_linear}. El operador efectivo $\sqrt{w}\,D$ se emplea
también para los grados de libertad y las barras de error.

\section{Verosimilitud Poisson en la distribución}
En modo Poisson el ajuste de histograma se repite con reponderado IRLS de la
$\sigma$ a partir del modelo (\S\ref{sec:weighting}), cerrando el lazo a través de
\code{lsq\_linear} hasta convergencia.

\section{Grados de libertad efectivos}\label{sec:effdof}

La regularización reduce los grados de libertad por debajo de $n$. El número
efectivo es la traza de la matriz sombrero:
\begin{equation}
  p_{\text{eff}} \;=\; \tr A(\alpha),\qquad
  A(\alpha)=X\bigl(X^\top W X+\alpha L^\top L\bigr)^{-1}X^\top W,
  \label{eq:effdof}
\end{equation}
calculable como $\tr\!\bigl[(X^\top WX+\alpha L^\top L)^{-1}X^\top WX\bigr]$. Se usa
en $\chi^2_\nu=\chi^2/(N-p_{\text{eff}})$, evitando el sesgo de contar los $n$
bins como parámetros independientes. Es invariante bajo el reescalado de
columnas de \S\ref{sec:cond}.

\section{Selección del parámetro de regularización}\label{sec:alpha}

Sobre una malla logarítmica de $\alpha$ se calculan tres criterios.

\paragraph{Curva L (máxima curvatura).} Con
$\rho(\alpha)=\log\lVert W^{1/2}(X\hat z-y)\rVert$ y
$\eta(\alpha)=\log\lVert L\hat{\bm P}\rVert$, la esquina maximiza
\begin{equation}
  \kappa(\alpha)=\frac{\rho'\eta''-\eta'\rho''}{\bigl((\rho')^2+(\eta')^2\bigr)^{3/2}}.
\end{equation}

\paragraph{Compromiso.} Mínima distancia normalizada al origen en el plano
$(\eta,\rho)$.

\paragraph{Validación cruzada generalizada (GCV).}
\begin{equation}
  \mathrm{GCV}(\alpha)=\frac{\lVert W^{1/2}(\bm y - X\hat z)\rVert^2}{\bigl(N-p_{\text{eff}}(\alpha)\bigr)^2},
  \qquad \hat\alpha=\argmin_\alpha \mathrm{GCV}(\alpha),
\end{equation}
que no requiere conocer la escala del ruido y es óptima en pérdida de predicción.

\section{Barras de error de \texorpdfstring{$P(\BHF)$}{P(BHF)}}\label{sec:perr}

Para el estimador regularizado $\hat{\bm z}=H^{-1}X^\top W\bm y$ con
$H=X^\top WX+\alpha L^\top L$, la covarianza linealizada es
\begin{equation}
  \widehat{\mathrm{Cov}}(\hat{\bm z}) \;=\; H^{-1}\,(X^\top WX)\,H^{-1}\cdot\frac{\chi^2}{N-p_{\text{eff}}},
\end{equation}
escalada por el $\chi^2_\nu$ efectivo (igual convención que el ajuste discreto). La
raíz de la diagonal correspondiente a los bins da $\sigma_{P_k}$. La banda
$P_k\pm\sigma_{P_k}$ se recorta a $P\ge0$ al representarla (la no-negatividad la
hace asimétrica en bins clavados a cero). En modo TV se usa el operador efectivo
$\sqrt{w}\,D$ en lugar de $L$.

\section{Distribuciones paramétricas}
Como alternativa al histograma regularizado, se ajustan formas cerradas con pocos
parámetros (mínimos cuadrados no lineales):
\begin{align}
  \text{Gaussiana:}\quad & P_k\propto \exp\!\bigl(-\tfrac12((c_k-\mu)/\varsigma)^2\bigr), & (b,s,A,\mu,\log\varsigma),\\
  \text{Binomial:}\quad & P_k\propto \binom{n-1}{k-1}p^{k-1}(1-p)^{n-k}, & (b,s,A,\mathrm{logit}\,p),\\
  \text{Fija:}\quad & P_k=\phi_k\ \text{(precargada)}, & (b,s,A).
\end{align}

\section{Refinamiento global (VARPRO)}\label{sec:refine}
Cuando se activa, un bucle externo no lineal optimiza los parámetros de forma
$\bm\eta=(\delta_g,\log\gamma_g,\dots)$; en cada evaluación se resuelve
\emph{íntegramente} el problema lineal interno \eqref{eq:dist-min}, devolviendo el
residuo al optimizador (esquema separable tipo VARPRO):
\begin{equation}
  \min_{\bm\eta}\;\bigl\|W^{1/2}\bigl[X(\bm\eta)\,\hat{\bm z}(\bm\eta)-\bm y\bigr]\bigr\|^2,
  \qquad \hat{\bm z}(\bm\eta)=\argmin_{\bm z}\eqref{eq:dist-min}.
\end{equation}

%======================================================================
\appendix
\part{Apéndices}

\section{Calibración del sextete}\label{ap:calib}
Las posiciones de referencia a $B_0=33\,$T se obtienen de
$r^{(0)}=\tfrac12[-10.657,-6.167,-1.677,1.677,6.167,10.657]\cdot(B_0/32.95)\,$mm/s.
De ellas se derivan los desdoblamientos Zeeman usados en \eqref{eq:Hkundig}:
\begin{equation}
  \omega_e^{(0)} = r^{(0)}_2-r^{(0)}_1 \approx 2{,}248\ \text{mm/s},\qquad
  \omega_g^{(0)} = -\bigl[(r^{(0)}_4-r^{(0)}_3)+\omega_e^{(0)}\bigr] \approx -3{,}928\ \text{mm/s}.
\end{equation}

\section{Cotas de los parámetros}\label{ap:bounds}
\begin{center}
\begin{tabular}{@{}lll@{}}
\toprule
Parámetro & Cota & Unidad\\
\midrule
$b$ (línea base) & $[0{,}7,\ 1{,}3]$ & --\\
$s$ (pendiente) & $[-5\cdot10^{-3},\ 5\cdot10^{-3}]$ & mm$^{-1}$s\\
$\delta$ & $[-2,\ 3]$ & mm/s\\
$\dEQ$ & $[0,\ 4]$ & mm/s\\
$\BHF$ & $[0,\ 50]$ & T\\
$\gamma_1$ & $[0{,}03,\ 2]$ & mm/s\\
$\gamma_2,\gamma_3$ & $[0{,}2,\ 3]$ & (relativa)\\
$d$ (profundidad) & $[0,\ 0{,}30]$ & --\\
$i_1,i_2,i_3$ & $[0,9]/[0,6]/[0,3]$ & (relativa)\\
$t$ (textura) & $[0,\ 1]$ & --\\
$\beta$ & $[0,\ 90]$ & grados\\
\bottomrule
\end{tabular}
\end{center}

\section{Función de Faddeeva}
$w(z)=e^{-z^2}\erfc(-iz)$, evaluada por \code{scipy.special.wofz}. La parte real
da el perfil de Voigt \eqref{eq:voigt}; para $z=iy$ con $y\in\mathbb R$,
$w(iy)=e^{y^2}\erfc(y)$ es real, lo que permite la normalización analítica al pico.

\section{Resumen de la pila de optimización}
\begin{center}
\begin{tabular}{@{}lll@{}}
\toprule
\textbf{Modo} & \textbf{Funcional} & \textbf{Solver}\\
\midrule
Discreto & $\tfrac12\sum_i\rho(r_i^2)$ & \code{least\_squares} (TRF) [+ DE opcional]\\
Distribución Tikhonov & $\|W^{1/2}(X\bm z-\bm y)\|^2+\alpha\|L\bm P\|^2$ & \code{lsq\_linear} (acotado)\\
Distribución TV & $\|W^{1/2}(X\bm z-\bm y)\|^2+\alpha\|D\bm P\|_1$ & IRLS sobre \code{lsq\_linear}\\
Distribución paramétrica & $\|W^{1/2}(X(\bm\theta)\bm z-\bm y)\|^2$ & \code{least\_squares} (NL)\\
Refinamiento global & $\|W^{1/2}(X(\bm\eta)\hat{\bm z}(\bm\eta)-\bm y)\|^2$ & externo NL + interno acotado\\
\bottomrule
\end{tabular}
\end{center}

\end{document}
